\subsection{An independent interval certificate}\label{push:section:E.4}
\label{sec:e4}

We write $\eta_D$ for the Dedekind eta function and set
\[
 h=z-(1-\tau)/2,\qquad
 \mathcal N^2(z,\tau)=
 \left|\frac{\vartheta_3(h,\tau)}{\eta_D(\tau)}\right|^2
 \exp\left(-\frac{2\pi(\Im h)^2}{\Im\tau}\right).
\]
The half-period theta identity identifies this with the norm used in
\cref{fam:src:C.1.1,fam:src:C.2.3}. The coefficient of $\Im c$ in $D_c$ is one. Consequently the
identity from the period calculation above which we use is
\[
 D_{\rm part}=\frac6\pi
 \log\left(\frac{\mathcal N}{2\sqrt3}\right).
\]
It follows directly by taking the imaginary part of the theta formula
for $\beta^{\rm part}$; its Gaussian factor removes
$6(\Im\mu)^2/\Im\tau$. Thus no numerical normalization is inferred
from the new computation. The analytic formula and the first
certificate are recorded in the period calculation above, \cref{fam:src:C.1.1,fam:src:C.2.3}. The result below supplies independent
arithmetic evidence for that same conclusion, with the independent certificate below.

\begin{theorem}[independent certification of every box]\label{push:src:E.4.1}
\StatementRecord{push:src:E.4.1}{proved}{Historical events: Claude Fable 5.2 (Anthropic), effort high, 16 November 2026: carried as background without fault found; used and rechecked in this run. Claude Fable 5.2 (Anthropic), effort high, 17 November 2026: carried as background without fault found. GPT-6 Astra (OpenAI), ChatGPT Pro mode, effort Pro, S6\_09\_PAPERS, 17 November 2026: targeted proof check. Also checked in this writing. Pre-continuation writing events: GPT-6 Astra (OpenAI), ChatGPT Pro mode, effort Pro, S6\_10\_PAPERS, 18 November 2026: computation or formal certificate; targeted proof check; used and rechecked in this run. Continuation metadata events: GPT-6 Astra (OpenAI), ChatGPT Pro mode, effort Pro, S6\_10\_PAPERS continuation, 19 November 2026: used and rechecked in this run; status-vocabulary reconciliation; metadata reconciliation only; no new mathematical check. \emph{Pre-continuation recorded proof scope.} The earlier attribution of code/I/replay\_interval\_leaves.output.json and replay\_family/f\_interval\_verify.stdout to this independent Arb theorem is withdrawn. Those executions used the earlier mpmath implementation and remain recorded under fam:src:C.2.3. The fresh unchanged Arb execution is separately recorded below. Read Johansson2016 Section2 pp2–4 and complete Section3 pp4–6, and official python-flint0.9.0 constructor/lower/upper/man\_exp/str semantics. These source-reading scopes are retained. The earlier reference to an already replayed Arb certificate was premature: that earlier execution used mpmath. The actual Arb implementation match and execution are recorded in the subsequent events. Read and rederived the full six-step proof and both complete original programs. Checked rational-to-Arb enclosures, exact dyadic Newton choices, the complex sup-norm contraction, finite pole-free Fourier enclosures, both moment tails, theta and logarithmic eta tails, upper-endpoint acceptance and the exact rational sweep. Wrote the explicit half-period theta identity and period/evenness transformations proving that the one-root norm enclosure controls both zeros. The successful contraction forces the selected preconditioner to be nonzero. The proof review is separate from the executable certificates. PASS: executed the original unchanged e4\_arb\_certificate.py with python-flint0.9.0,192bits,all19701boxes,four workers. All19701 strict root contractions and norm bounds succeeded,0failures. The largest upper endpoint is enclosed by1.999877571767410524070536478244687423578792301873142913 plus or minus9.26e-56,below1.999877571767410525; worst index17866. The input SHA256 is87b7c2dc672e98cab926c092ebdc757373a6817bb651c3472235a42956b5c9ab. Only rational endpoints from the earlier certificate are imported. Full outputs and identical source hashes are in code/I/replay\_arb and analysis/I\_arb\_replay\_manifest.json. PASS: executed the original unchanged e4\_cover\_and\_local.py. The exact Fraction sweep checks19911leaves,438vertical edges,area23/40,disjoint interiors and complete coverage;36local rectangles lie in the prescribed disc and174excluded rectangles lie below the unit circle. The separately implemented256bit Arb local-constant inequalities also all passed. Execution of those constants grants no additional whole-proof review of the older local argument.}{the hypotheses, notation and definitions stated here; \citet[Sections 2--3]{Johansson2016}; \citet{PythonFlint2026}}



For all 19701 rational rectangles in the supplied certificate, and for
each $\tau$ in the rectangle, both zeros of $\wp_\tau$ satisfy
$\mathcal N^2<2$. The new largest certified upper bound is less than
\[
 1.999877571767410525.
\]
The 19701 rectangles, 36 local rectangles and 174 excluded rectangles
have disjoint interiors and cover exactly
$[0,1/2]\times[17/20,2]$. Every local rectangle is contained in
$|\tau-i|<3/20000$, and every excluded rectangle satisfies
$|\tau|<1$ throughout.
\end{theorem}



\begin{proof}

\begin{enumerate}[label=\arabic*.,leftmargin=*]
\item The programs \path{code/I/replay_arb/S6_02_PUSH/code/e4_arb_certificate.py} and
\path{code/I/replay_arb/S6_02_PUSH/code/e4_cover_and_local.py} use respectively 192-bit and 256-bit
Arb arithmetic through python-flint 0.9.0. The arithmetic implementation
is independent of mpmath. We import neither mpmath nor any routine
from the previous certificate. The only numerical data taken from that
certificate are its rational rectangle endpoints. We recompute roots,
preconditioners, radii, contraction bounds, function values and tails.
Arb's midpoint-radius representation and its outward error bounds are
described by \citet[Section~2, pp.~2--4, and Section~3, pp.~4--6]{Johansson2016}; the precise
constructor and endpoint semantics used here are documented in
\citet{PythonFlint2026}. The final comparison is made with an upper
endpoint, before conversion to a decimal string.

\item For a rectangle $Q$, let $z_0$ and $A\ne0$ be the exact dyadic complex
numbers selected by the midpoint Newton iteration. That iteration is
only a choice of constants. For $Z=z_0+[-r,r]+i[-r,r]$ we independently
enclose $\wp_\tau(z_0)$ and $\wp'_\tau(Z)$ for $\tau\in Q$.
We accept only if
\[
 \begin{split}
 &\left|-A\wp_\tau(z_0)
 +(1-A\wp'_\tau(Z))(Z-z_0)\right|_\infty<r,\\
 &|\Re(1-A\wp'_\tau(Z))|
       +|\Im(1-A\wp'_\tau(Z))|<1.
 \end{split}
\]
Here absolute values on enclosures mean their outward upper bounds.
The derivative enclosure and integration along a line segment show
that $z\mapsto z-A\wp_\tau(z)$ is a strict contraction of $Z$ into
itself, uniformly for $\tau\in Q$. Thus it has a genuine zero in $Z$.
The finite Fourier enclosures exclude the poles on $Z$, and the strict
Lipschitz inequality forces $A\ne0$. The two zeros of $\wp$ are
negatives modulo the period lattice, counted with multiplicity. The
exact half-period transformation, obtained by replacing the summation
index in the theta series, is
\[
 \vartheta_3\bigl(z-(1-\tau)/2,\tau\bigr)
   =i e^{-\pi i\tau/4-\pi i z}\vartheta_1(z,\tau).
\]
Writing $v=\Im z$ and $y=\Im\tau$, its modulus and Gaussian factor
cancel to give
\[
 \mathcal N(z,\tau)=
 \left|\frac{\vartheta_1(z,\tau)}{\eta_D(\tau)}\right|
 e^{-\pi v^2/y}.
\]
Now $\vartheta_1(-z)=-\vartheta_1(z)$,
$\vartheta_1(z+1)=-\vartheta_1(z)$ and
$\vartheta_1(z+\tau)=-e^{-\pi i\tau-2\pi i z}\vartheta_1(z)$.
In the last transformation the modulus contributes
$e^{\pi y+2\pi v}$, while the Gaussian contributes its reciprocal.
Thus the displayed norm is even and lattice-periodic, proving directly
that the one-root enclosure controls both zeros. No continuity or
success of Newton iteration is used in that conclusion.

\item We use the normally convergent cosine expansion
\[
 \wp_\tau(z)=\frac{\pi^2}{\sin^2\pi z}-\frac{\pi^2}{3}
 +8\pi^2\sum_{n\geq1}\frac{nq^n}{1-q^n}
                      (1-\cos 2\pi nz),\qquad q=e^{2\pi i\tau},
\]
and differentiate it in $z$ for $\wp'$. This identity follows by twice
differentiating the normally convergent Jacobi product and inserting
$E_2=1-24\sum_{n\geq1}nq^n/(1-q^n)$; this also verifies the
Fourier formula used in the period calculation above. The new program sums
32 terms directly in cosine form. Put $r_q\geq|q|$,
$H_z\geq|\Im z|$, and $R_q=r_q e^{2\pi H_z}<1$.
For $j=1,2$, the inequality $(N+k)^j\leq(N+1)^jk^j$ gives
\[
 \begin{aligned}
 U_1(a)&=\frac{(N+1)a^{N+1}}{(1-a)^2},\\
 U_2(a)&=\frac{(N+1)^2a^{N+1}(1+a)}{(1-a)^3},
 \qquad
 \sum_{n>N}n^j a^n\leq U_j(a).
 \end{aligned}
\]
We therefore add complex rectangles enclosing discs of radii
\[
 \frac{8\pi^2}{1-r_q}\bigl(U_1(r_q)+U_1(R_q)\bigr),
 \qquad \frac{16\pi^3}{1-r_q}U_2(R_q)
\]
to the two partial sums. These estimates follow from
$|\cos z|,|\sin z|\leq e^{|\Im z|}$ and
$|1-q^n|\geq1-r_q$. They differ from the moment formulas used in
the old implementation.

\item For $\vartheta_3$ we sum
$1+2\sum_{n=1}^{12}e^{\pi i n^2\tau}\cos(2\pi nh)$.
With $T_*\leq\Im\tau$ and $H_*\geq|\Im h|$, the remaining
absolute value is at most
\[
 \frac{2e^{-169\pi T_*+26\pi H_*}}
 {1-e^{-27\pi T_*+2\pi H_*}}.
\]
The ratio in this expression is tested to be less than one. For eta
we independently sum its logarithm,
\[
 \log\eta_D(\tau)=\frac{\pi i\tau}{12}
          +\sum_{n=1}^{24}\log(1-q^n)+\mathcal E,
 \qquad
 |\mathcal E|\leq
 \frac{r_q^{25}}{(1-r_q)(1-r_q^{25})}.
\]
The logarithms are holomorphic on $|q^n|<1$. The tail follows by
summing $|\log(1-z)|\leq|z|/(1-|z|)$.
We enclose $\mathcal N^2$ as
\[
 ((\Re\vartheta_3)^2+(\Im\vartheta_3)^2)
 \exp\left(-2\Re\log\eta_D-
                   \frac{2\pi(\Im h)^2}{\Im\tau}\right),
\]
and test its upper endpoint against the exact integer 2.

\item All 19701 tests succeeded. The full output records each new enclosure;
the largest is enclosed by
\[
 1.999877571767410524070536478244687423578792301873142913
       \ \pm 9.26\cdot10^{-56}.
\]
This is the enclosure of a certified upper bound, not an estimate of
the actual maximum. Its box has input index 17866, with zero-based
indexing. The file digest of the input is recorded in the output.

\item The cover check is exact rational arithmetic. At all 438 distinct
vertical edges we maintain the active closed vertical intervals. In
every open vertical slab they form an ordered partition of
$[17/20,2]$, each with multiplicity one. The last active set is empty;
all leaves lie in the root rectangle, and the sum of their exact areas
is $23/40$. The closed rectangles also cover their common edges.
The omitted and excluded conditions are checked separately using
rational endpoint comparisons. Finally,
$2(1/10000)^2<(3/20000)^2$ puts the entire omitted square inside the
analytic disc used below. This verifies both the inequalities and
the coverage, rather than just replaying the old list of successful
flags.
\end{enumerate}

\end{proof}

\begin{theorem}[quantitative local maximum]\label{push:src:E.4.2}
\StatementRecord{push:src:E.4.2}{proved}{Historical events: Claude Fable 5.2 (Anthropic), effort high, 16 November 2026: carried as background without fault found. GPT-6 Astra (OpenAI), ChatGPT Pro mode, effort Pro, S6\_09\_PAPERS, 17 November 2026: targeted proof check. Also checked in this writing. Pre-continuation writing events: GPT-6 Astra (OpenAI), ChatGPT Pro mode, effort Pro, S6\_10\_PAPERS, 18 November 2026: targeted proof check. Continuation metadata events: GPT-6 Astra (OpenAI), ChatGPT Pro mode, effort Pro, S6\_10\_PAPERS continuation, 19 November 2026: used and rechecked in this run; status-vocabulary reconciliation; metadata reconciliation only; no new mathematical check. \emph{Pre-continuation recorded proof scope.} Independently re-proved the quantitative local maximum through the explicit sigma-to-Jacobi norm identity and the checked global maximum; the small root-disc contains only the i representative. Earlier constants and interval proof remain in the compendium as a second proof.}{\cref{ari:C.1.1,ari:C.8.15}}

If $|\tau-i|\leq3/20000$ and $\wp_\tau(z)=0$, then
$\mathcal N(z,\tau)\leq\sqrt2$, with equality precisely at
$\tau=i$. The explicit numerical second proof is retained in the compendium;
the global norm theorem gives the following short proof.
\end{theorem}

\begin{proof}
\begin{enumerate}[label=\arabic*.,leftmargin=*]
\item For $\Lambda=\Z+\Z\tau$, the normalization identities are
\[
 \Delta_\Lambda=(2\pi)^{12}\eta_D(\tau)^{24},\qquad
 \sigma_\Lambda(z)=
 \frac{\vartheta_1(z,\tau)}{2\pi\eta_D(\tau)^3}
                 e^{\eta_1z^2/2}.
\]
The real-linear quasiperiod formula in \ref{ari:C.1.1} then gives
\[
 |\Delta_\Lambda|^{1/12}|\sigma_\Lambda(z)|
 e^{-\Re(z\eta(z))/2}
 =\left|\frac{\vartheta_1(z,\tau)}{\eta_D(\tau)}\right|
            e^{-\pi(\Im z)^2/\Im\tau}
 =\mathcal N(z,\tau).
\]
For $\wp_\tau(z)=0$, Theorem~\ref{ari:C.8.15} gives the
global maximum $\log\sqrt2$ of its logarithm, exactly on the
orbit of $i$. In the disc $|\tau-i|\le3/20000$ the only such
representative is $i$: a modular image with $c^2+d^2\ge2$
has imaginary part at most $1/2$, and the remaining images
are integral translates of $i$. Exponentiating proves the
stated inequality and its equality locus.
\end{enumerate}
\end{proof}




\clearpage\appendix
